\newpage
\section{CPA Equation of State}
\subsection{Compressibility Factor}
The density of the mixture can be obtained from the compressibility factor, which can be derived from the residual Helmholtz free energy:
\begin{equation}
     Z = 1 - \frac{v}{RT}\left (  \frac{\partial A^r}{\partial V} \right )_{T,N_i}.
    \label{Eq:Z}
\end{equation}

The compressibility factor can be divided into terms representing different interactions. We also add a term related to the change in the relative permittivity (perm), which is from the DH and Born term (more details on the origin of this term are presented afterward):
\begin{equation}
        Z = Z^\text{phys} + Z^\text{assoc}  + Z^\text{DH} + Z^\text{Born} + Z^\text{Perm}.
        \label{Eq:Z_eCPA}
\end{equation}

The compressibility factor expression for the physical term is given by:
\begin{equation}
        Z^\text{phys} = \frac{Z}{Z-B} - \frac{A}{Z+B},
        \label{Eq:Z_phys}
\end{equation}
where $A$ and $B$ are the dimensionless SRK parameters:
\begin{equation}
    A=\frac{aP}{R^2T^2},
    \label{Eq:CPA_A}
\end{equation}
\begin{equation}
    B=\frac{bP}{RT}.
    \label{Eq:CPA_B}
\end{equation}

The association expression for the compressibility factor is given by:
\begin{equation}
        Z^\text{assoc} =  \frac{1}{2}\left ( 1+\frac{\eta }{g_\eta}\frac{\mathrm{d} g_\eta}{\mathrm{d} \eta} \right )\sum_i x_i \sum_{A \in i}  \left ( \chi_{A_i} -1 \right ).
        \label{Eq:Z_assoc}
\end{equation}

In the DH and Born terms, we split them into a contribution at fixed relative permittivity and varying relative permittivity, which originates $Z^\text{Perm}$:
\begin{equation}
        - \frac{v}{RT}\left (  \frac{\partial A^\text{DH}}{\partial V} \right )_{T,N_i} = - \frac{v}{RT}\left (  \frac{\partial A^\text{DH}}{\partial V} \right )_{T,N_i,\epsilon_r} - \frac{v}{RT}\left (  \frac{\partial A^\text{DH}}{\partial \epsilon_r} \right )_{T, V, N_i}\left (  \frac{\partial \epsilon_r}{\partial V} \right )_{T, N_i},
        \label{Eq:ZDH_cont}
\end{equation}
\begin{equation}
        Z^\text{DH} = - \frac{v}{RT}\left (  \frac{\partial A^\text{DH}}{\partial V} \right )_{T,N_i,\epsilon_r} = \frac{v}{4\pi N_A  \sum_j x_j z_j^2}\sum_j x_j z_j^2 \left [ X_j - \frac{\kappa^3}{2\left ( 1 + \kappa \sigma_j  \right )} \right ],
        \label{Eq:ZDH}
\end{equation}
\begin{equation}
        - \frac{v}{RT}\left (  \frac{\partial A^\text{Born}}{\partial V} \right )_{T,N_i} = - \frac{v}{RT}\left (  \frac{\partial A^\text{Born}}{\partial V} \right )_{T,N_i,\epsilon_r} - \frac{v}{RT}\left (  \frac{\partial A^\text{Born}}{\partial \epsilon_r} \right )_{T, V, N_i}\left (  \frac{\partial \epsilon_r}{\partial V} \right )_{T, N_i},
        \label{Eq:ZBorn_cont}
\end{equation}
\begin{equation}
        Z^\text{Born}  = - \frac{v}{RT}\left (  \frac{\partial A^\text{Born}}{\partial V} \right )_{T,N_i,\epsilon_r} = 0,
        \label{Eq:ZBorn}
\end{equation}
\begin{equation}
        Z^\text{Perm} = - \frac{v}{RT}\left [ \left (  \frac{\partial A^\text{DH}}{\partial \epsilon_r} \right )_{T, V, N_i}  +\left (  \frac{\partial A^\text{Born}}{\partial \epsilon_r} \right )_{T, V, N_i} \right ]\left (  \frac{\partial \epsilon_r}{\partial V} \right )_{T, N_i},
        \label{Eq:ZPerm_cont}
\end{equation}
\begin{equation}
        Z^{\text{Perm}} =  -\left [ \frac{\kappa^2}{8\pi N_A c \epsilon_r \sum_j x_j z_j^2}\sum_j \left (  \frac{\kappa x_j z_j^2}{1+\kappa \sigma_j} - \frac{ x_j z_j^2}{R_{B,j}} \right )  \right ] V \left  (  \frac{\partial \epsilon_r}{\partial V} \right )_{T, N_i}.
        \label{Eq:ZPerm}
\end{equation} 

We, therefore, need an expression for the variation of the dielectric constant with volume.

\subsection{Fugacity Coefficient}
When two phases are in equilibrium, temperature, pressure, and fugacity of each component in both phases are equal. In terms of chemical equilibrium, the fugacity equality of component $i$ between phases I and II is given by:
\begin{equation}
     x_i^{I}\phi_i^I = x_i^{II}\phi_i^{II},
    \label{Eq:xPHIxPHI}
\end{equation}
where $\phi_i$ is the fugacity coefficient of component $i$, which can be derived from the residual Helmholtz free energy as:
\begin{equation}
     \ln{\phi_i} =   \frac{1}{RT}\left (  \frac{\partial A^r}{\partial N_i} \right )_{T,V,N_j} - \ln{Z}.
    \label{Eq:lnPHI}
\end{equation} 

Similarly to the compressibility factor, the fugacity coefficient can be divided into terms representing each one of the interactions:
\begin{equation}
    \ln{\phi_i} = \ln{\phi_i^\text{phys}} + \ln{\phi_i^\text{assoc}}  + \ln{\phi_i^\text{DH}}  + \ln{\phi_i^\text{Born}}  + \ln{\phi_i^\text{Perm}}.
    \label{Eq:eCPA_phi}
\end{equation}

Separating the relative permittivity contribution as done in the compressibility factor, the expressions for the physical, association, DH and Born and relative permittivity fugacity coefficients are given, respectively, by:
\begin{equation}
\begin{aligned}
\ln{\phi_i^{\text{phys}}} = &  -\ln{\left ( Z-B \right )}+\frac{B_i}{B}\left ( \frac{B}{Z-B}-\frac{A}{Z+B} \right ) \  \\
      &   -\left \{ \frac{A_i}{B_i}+\frac{1}{B\ln{2}}\left [ B_i\frac{g^{E\infty}}{RT} - \sum_jx_j\left ( B_i\frac{\Delta U_{ij}}{RT} + B_j\frac{\Delta U_{ji}}{RT} \right ) \right ] \right \}\ln{\left (1+\frac{B}{Z}  \right )},
      \label{Eq:eCPA_phi_phys}
\end{aligned}
\end{equation}
\begin{equation}
    \ln{\phi_i^{\text{assoc}}}  = \sum_{A \in i} \eta_{A_i}\ln\chi_{A_i}+\frac{B_i}{8g_{\eta}Z}\frac{\mathrm{d} g_\eta}{\mathrm{d} \eta} \sum_j x_j  \sum_{A \in j}\left (\chi_{A_j}-1\right ),
    \label{Eq:CPA_phi_assoc}
\end{equation}
\begin{equation}
    \ln{\phi_i^{\text{DH}}} = \frac{z_i^2v}{4\pi N_A  \sum_j x_j z_j^2} \left [  \frac{\sum_j x_j z_j^2 X_j}{\sum_j x_j z_j^2} - X_i - \frac{\kappa^3}{2 \sum_j x_j z_j^2} \sum_j \left (  \frac{x_jz_j^2}{1+\kappa\sigma_j}\right ) \right ],
    \label{Eq:Phi_DH}
\end{equation}
\begin{equation}
    \ln{\phi_i^{\text{Born}}}  =   \frac{N_A e^2 z_i^2}{8\pi \epsilon_0RT R_{B,i} } \left ( \frac{1}{\epsilon_r} - 1\right ),
    \label{Eq:Phi_Born}
\end{equation}
\begin{equation}
        \ln{\phi_i^{\text{Perm}}} =  \left [ \frac{\kappa^2}{8\pi N_A c \epsilon_r \sum_j x_j z_j^2}\sum_j \left (  \frac{\kappa x_j z_j^2}{1+\kappa \sigma_j} - \frac{ x_j z_j^2}{R_{B,j}} \right )  \right ] N \left  (  \frac{\partial \epsilon_r}{\partial N_i} \right )_{T, V, N_j}.
        \label{Eq:PhiPerm}
\end{equation} 

Here, we need an expression for the variation of the dielectric constant with composition.

\subsection{Relative Permitivitty}
To determine the relative permittivity contribution to the compressibility factor (Eq. \ref{Eq:ZPerm}) and to the fugacity coefficient (Eq. \ref{Eq:PhiPerm}), one needs to evaluate the derivatives $\partial \epsilon_r / \partial V$ and $\partial \epsilon_r / \partial N_i$. We show how to obtain these derivatives from eCPA Eos and the relative permittivity primitive model from \citeauthor{Maribo2013a} \cite{Maribo2013}.  Assuming water as the only polar molecule in the mixture with a net dipole moment, we can define function $F\left (  \epsilon_r, \epsilon_\infty, N_i, V \right )$:
\begin{equation}
        F\left (  \epsilon_r, \epsilon_\infty, N_i, V \right ) = \frac{\left (2\epsilon_r + \epsilon_\infty \right ) \left ( \epsilon_r - \epsilon_\infty \right )}{\epsilon_r\left ( \epsilon_\infty + 2 \right )^2} - \frac{N_A}{9\epsilon_0k_BTv}x_wg_w\mu_{0w}^2 =0.
        \label{Eqeps_simp}
\end{equation} 

The water Kirkwood g-factor and the probability of association between water and other molecules are given by \cite{Maribo2013}: 
\begin{equation}
        g_w = 1 + \frac{z_{ww}\mathcal{P}_{ww}\cos{\gamma_{ww}}}{\mathcal{P}_{w} \cos{\theta_{ww}+1}},
        \label{Kirkwood_simp}
\end{equation} 
\begin{equation}
        \mathcal{P}_{w} = \sum_j \mathcal{P}_{wj},
        \label{Prob_association_total}
\end{equation} 
\begin{equation}
        \mathcal{P}_{wj} = \eta_{A_w}c x_j\chi_{A_w} \chi_{B_j} \delta_{A_iB_j},
        \label{P_wj}
\end{equation} 
where $\eta_{A_w}$ is the number of associating sites type A in water (2 for the 4C scheme).

To obtain the derivatives $\partial \epsilon_r / \partial V$ and $\partial \epsilon_r / \partial N_i$, we differentiate the function $F$: 
\begin{equation}
        \mathrm{d} F = \left ( \frac{\partial F}{\partial \epsilon_r} \right )\mathrm{d}\epsilon_r + \left ( \frac{\partial F}{\partial \epsilon_\infty} \right )\mathrm{d}\epsilon_\infty + \left ( \frac{\partial F}{\partial V} \right )\mathrm{d}V + \sum_i\left ( \frac{\partial F}{\partial N_i} \right )\mathrm{d}N_i = 0.
        \label{Ffunction}
\end{equation} 

To simplify the notation, we omit each term that is kept constant at each partial derivative. The dielectric constant derivatives can be expressed as:
\begin{equation}
        V\left ( \frac{\partial \epsilon_r}{\partial V} \right )_{T,Ni} = -\left ( \frac{\partial F}{\partial \epsilon_r} \right )^{-1}\left [  \left ( \frac{\partial F}{\partial \epsilon_\infty} \right )V\left ( \frac{\partial \epsilon_\infty}{\partial V} \right ) + V\left ( \frac{\partial F}{\partial V} \right ) \right ],
        \label{Ffunction_V}
\end{equation} 
\begin{equation}
        N\left ( \frac{\partial \varepsilon_r}{\partial N_i} \right )_{T, V, Nj} = -\left ( \frac{\partial F}{\partial \varepsilon_r} \right )^{-1}\left [  \left ( \frac{\partial F}{\partial \varepsilon_\infty} \right )N\left ( \frac{\partial \varepsilon_\infty}{\partial N_i} \right ) + N\left ( \frac{\partial F}{\partial N_i} \right ) \right ].
        \label{Ffunction_N}
\end{equation} 

Each term of Eqs. \eqref{Ffunction_V} and \eqref{Ffunction_N} is given by the following expressions:
\begin{equation}
         \left ( \frac{\partial F}{\partial \epsilon_r} \right )=\frac{2\epsilon_r^2+\epsilon_\infty^2}{\epsilon_r^2\left ( \epsilon_\infty+2 \right )^2},
        \label{dFder}
\end{equation}
\begin{equation}
          \left ( \frac{\partial F}{\partial \epsilon_\infty} \right )=\frac{\epsilon_r\epsilon_\infty - 4 \epsilon_\infty -4\epsilon_r^2-2\epsilon_r}{\epsilon_r\left ( \epsilon_\infty+2 \right )^3},
        \label{dFdeinf}
\end{equation}
\begin{equation}
          V\left ( \frac{\partial \epsilon_\infty}{\partial V} \right )=- \frac{3M_{\epsilon_\infty}}{\left ( 1 - M_{\epsilon_\infty} \right )^2},
        \label{deinfdV}
\end{equation}
\begin{equation}
          N\left ( \frac{\partial \epsilon_\infty}{\partial N_i} \right )=\frac{N_A c \alpha_{0i}}{\epsilon_0 \left ( 1 - M_{\epsilon_\infty} \right )^2},
        \label{deinfdNi}
\end{equation}
\begin{equation}
           V\left ( \frac{\partial F}{\partial V} \right )= V \left ( \frac{\partial F}{\partial V} \right )_{g_w} + \left ( \frac{\partial F}{\partial g_w} \right ) V\left ( \frac{\partial g_w}{\partial V} \right ) = \frac{N_Ac x_w\mu_{0w}^2}{9\epsilon_0k_BT}\left [ g_w -  V \left ( \frac{\partial g_w}{\partial V} \right ) \right ],
        \label{dFdV}
\end{equation}
\begin{equation}
            N\left ( \frac{\partial F}{\partial N_i} \right )= N \left ( \frac{\partial F}{\partial N_i} \right )_{g_w} + \left ( \frac{\partial F}{\partial g_w} \right ) N\left ( \frac{\partial g_w}{\partial N_i} \right ) = -\frac{N_Ac x_w \mu_{0w}^2}{9\varepsilon_0k_BT}\left [ \frac{g_w \Gamma_{iw}}{x_w} +  N \left ( \frac{\partial g_w}{\partial N_i} \right ) \right ],
        \label{dFdNi}
\end{equation}
where $\Gamma_{ij} = 1$ if $i = j$, and $\Gamma_{ij} = 0$ otherwise, and
\begin{equation}
        M_{\epsilon_\infty} = \frac{\epsilon_\infty-1}{\epsilon_\infty+2}.
        \label{Eq:eps_infty}
\end{equation} 

Now we evaluate two special cases. In the first one, water is the only associating molecule and we can obtain an analytical expression for $\partial g_w / \partial V$ and $\partial g_w / \partial N_i$. In the second case, besides water self-association, we also have \ce{CO2}-\ce{H2O} cross-association, in such a way that numerical methods are used to obtain the derivatives of $\partial \chi_w / \partial V$ and $\partial \chi_w / \partial N_i$

\subsubsection{Special Case I: Water as single associating molecule}
Assuming that $\varphi = V$, $N$ or $N$ if $\Phi = V$, $N_w$ or $N_c$.
\begin{equation}
           \varphi \left ( \frac{\partial g_w}{\partial \Phi} \right ) = \left ( \frac{\partial g_w}{\partial \mathcal{P}_{ww}} \right ) \left ( \frac{\partial \mathcal{P}_{ww}}{\partial \chi_{w}} \right )\varphi\left ( \frac{\partial \chi_w}{\partial \Phi} \right ),
        \label{dPhidgw}
\end{equation}
\begin{equation}
        \mathcal{P}_{w} = \mathcal{P}_{ww} = 1 - \chi_w,
        \label{P_w_onlywater}
\end{equation} 
\begin{equation}
         \left ( \frac{\partial g_w}{\partial \mathcal{P}_{ww}} \right ) = \frac{g_w-1}{\mathcal{P}_{ww} \left ( \mathcal{P}_{ww} \cos{\theta_{ww}}  + 1 \right )},
        \label{dgwdPw}
\end{equation} 
\begin{equation}
         \left ( \frac{\partial \mathcal{P}_{ww}}{\partial \chi_{w}} \right ) = -1,
        \label{dPwdchi}
\end{equation} 
\begin{equation}
         V\left ( \frac{\partial \chi_w}{\partial V} \right ) = \frac{1}{4c x_w\delta}\left ( \frac{1+4c x_w\delta}{\sqrt{1+8c x_w\delta}} -1\right )\left ( 1 + \frac{\eta }{g_\eta}\frac{\mathrm{d} g_\eta}{\mathrm{d} \eta}  \right ),
        \label{dchidV}
\end{equation} 
\begin{equation}
         N\left ( \frac{\partial \chi_w}{\partial N_i} \right ) = \frac{1}{4c x_w\delta}\left ( 1-\frac{1+4c x_w\delta}{\sqrt{1+8c x_w\delta}}\right )\left ( \frac{\Gamma_{iw}}{x_w} + \frac{c b_i }{4 g_\eta}\frac{\mathrm{d} g_\eta}{\mathrm{d} \eta}  \right ).
        \label{dchidNi}
\end{equation} 



\subsubsection{Special Case II: Water self-association + \ce{CO2} cross-association}

\begin{equation}
        \mathcal{P}_{w} =  \mathcal{P}_{ww} + \mathcal{P}_{wc},
        \label{Prob_association_water}
\end{equation} 
\begin{equation}
        \mathcal{P}_{ww} = 2c x_w\chi_{w}^2 \delta_w,
        \label{P_ww}
\end{equation} 
\begin{equation}
        \mathcal{P}_{wc} = 2c x_c\chi_{w}\chi_{c} S_{cw}\delta_w,
        \label{P_wc}
\end{equation} 
\begin{equation}
\begin{aligned}
           \varphi \left ( \frac{\partial g_w}{\partial \Phi} \right ) & =  \left ( \frac{\partial g_w}{\partial \mathcal{P}_{ww}} \right ) \varphi\left ( \frac{\partial \mathcal{P}_{ww}}{\partial \Phi} \right ) + \left ( \frac{\partial g_w}{\partial \mathcal{P}_{w}} \right ) \varphi\left ( \frac{\partial \mathcal{P}_{w}}{\partial \Phi} \right ) \\
           & = \frac{g_w-1 }{ \mathcal{P}_{ww} }  \varphi\left ( \frac{\partial \mathcal{P}_{ww}}{\partial \Phi} \right ) + \frac{\left (g_w-1 \right )\cos{\theta_{ww}}}{ \mathcal{P}_{w} \cos{\theta_{ww}}  + 1}  \left [  \varphi\left ( \frac{\partial \mathcal{P}_{ww}}{\partial \Phi} \right ) + \varphi\left ( \frac{\partial \mathcal{P}_{wc}}{\partial \Phi} \right ) \right ].
        \label{dgwdPhi}
\end{aligned}
\end{equation}

$\mathcal{P}_{ww} = f \left ( N_w, V, \delta_w, \chi_w\right )$, whereas $\mathcal{P}_{wc} = f \left ( N_c, V, \delta_w, \chi_w ,  \chi_c, \right )$. Their derivatives with respect of $V$ and $N_i$ are then given, respectively, by: 
\begin{equation}
        V \frac{\partial \mathcal{P}_{ww}}{\partial V} = \mathcal{P}_{ww} \left [ -1 + \frac{\delta_w \eta}{g_\eta} \frac{\mathrm{d} g_\eta} {\mathrm{d} \eta} + \frac{2}{\chi_w} V \left ( \frac{\partial \chi_w}{\partial V} \right )\right ],
        \label{dPwwdV}
\end{equation}
\begin{equation}
        V \frac{\partial \mathcal{P}_{wc}}{\partial V} = \mathcal{P}_{wc} \left [ -1 + \frac{\delta_w \eta}{g_\eta} \frac{\mathrm{d} g_\eta} {\mathrm{d} \eta} + \frac{1}{\chi_w} V \left( \frac{\partial \chi_w}{\partial N_i} \right ) + \frac{1}{\chi_c} V \left( \frac{\partial \chi_c}{\partial N_i} \right )\right ],
        \label{dPwcdV}
\end{equation}
\begin{equation}
        N \frac{\partial \mathcal{P}_{ww}}{\partial N_i} = \mathcal{P}_{ww} \left [ \frac{\Gamma_{iw}}{x_w} +  \frac{\delta_w \eta b_i}{g_\eta  b} \frac{\mathrm{d} g_\eta} {\mathrm{d} \eta} + \frac{2}{\chi_w} N \left( \frac{\partial \chi_w}{\partial N_i} \right )\right ],
        \label{dPwwdNi}
\end{equation}
\begin{equation}
        N \frac{\partial \mathcal{P}_{wc}}{\partial N_i} = \mathcal{P}_{wc} \left [\frac{\Gamma_{ic}}{x_c} + \frac{\delta_w \eta b_i}{g_\eta  b} \frac{\mathrm{d} g_\eta} {\mathrm{d} \eta} + \frac{1}{\chi_w} N \left( \frac{\partial \chi_w}{\partial N_i} \right ) +\frac{1}{\chi_c} N \left( \frac{\partial \chi_c}{\partial N_i} \right )\right ]
        \label{dPwcdNi}.
\end{equation}

The derivatives of $\chi_i$ with respect to $V$ and $N_i$are obtained numerically. First, we define two functions $G$ and $H$:
\begin{equation}
    G \left(\chi_w, \chi_c, \delta_w, V, N_w \right ) = \chi_c -  \frac{V}{V+2N_w\chi_w S_{cw}\delta_w} =0,
    \label{Gfunction}
\end{equation}
\begin{equation}
    H \left(\chi_w,  \chi_c, \delta_w, V, N_w, N_c \right ) = \chi_w -  \frac{V}{V+2N_w\chi_w\delta_w+2N_c\chi_cS_{cw}\delta_w}=0,
    \label{Hfunction}
\end{equation}
\begin{equation}
        \mathrm{d} G = \left ( \frac{\partial G}{\partial \chi_w} \right )\mathrm{d}\chi_w +\left ( \frac{\partial G}{\partial \chi_c} \right )\mathrm{d}\chi_c + \left ( \frac{\partial G}{\partial \delta_w} \right )\mathrm{d}\delta_w + \left ( \frac{\partial G}{\partial V} \right )\mathrm{d}V + \left ( \frac{\partial G}{\partial N_w} \right )\mathrm{d}N_w = 0,
        \label{Gfunction_derivative}
\end{equation} 
\begin{equation}
        \mathrm{d} H = \left ( \frac{\partial H}{\partial \chi_w} \right )\mathrm{d}\chi_w + \left ( \frac{\partial H}{\partial \chi_c} \right )\mathrm{d}\chi_c + \left ( \frac{\partial H}{\partial \delta_w} \right )\mathrm{d}\delta_w + \left ( \frac{\partial H}{\partial V} \right )\mathrm{d}V + \left ( \frac{\partial H}{\partial N_w} \right )\mathrm{d}N_w  + \left ( \frac{\partial H}{\partial N_c} \right )\mathrm{d}N_c = 0,
        \label{Hfunction_derivative}
\end{equation} 
\begin{equation}
        \varphi \left ( \frac{\partial \chi_c}{\partial \Phi} \right ) = -\left ( \frac{\partial G}{\partial \chi_c} \right )^{-1}\left [  \left ( \frac{\partial G}{\partial \chi_w} \right )\varphi\left ( \frac{\partial \chi_w}{\partial \Phi} \right ) + \left ( \frac{\partial G}{\partial \delta_w} \right )\varphi\left ( \frac{\partial \delta_w}{\partial \Phi} \right ) + \varphi\left ( \frac{\partial G}{\partial \Phi} \right ) \right ],
        \label{Gfunction_Phi}
\end{equation} 
\begin{equation}
        \varphi \left ( \frac{\partial \chi_w}{\partial \Phi} \right ) = -\left ( \frac{\partial H}{\partial \chi_w} \right )^{-1}\left [  \left ( \frac{\partial H}{\partial \chi_c} \right )\varphi\left ( \frac{\partial \chi_c}{\partial \Phi} \right ) + \left ( \frac{\partial H}{\partial \delta_w} \right )\varphi\left ( \frac{\partial \delta_w}{\partial \Phi} \right ) + \varphi\left ( \frac{\partial H}{\partial \Phi} \right ) \right ].
        \label{Hfunction_Phi}
\end{equation} 

We have three new unknowns: $\partial \chi_w / \partial V, \partial \chi_w / \partial N_w$, and  $\partial \chi_w / \partial N_c$ that can be solved numerically by coupling Eqs. \ref{Gfunction_Phi} and \ref{Hfunction_Phi}, knowing that
$\partial \chi_c / \partial N_c = 0$. All the other derivatives are given next:
\begin{equation}
        \left ( \frac{\partial G}{\partial \chi_w} \right ) = 2 x_w c S_{cw} \delta_w \chi_w^2,
        \label{dGdchiw}
\end{equation} 
\begin{equation}
        \left ( \frac{\partial G}{\partial \chi_c} \right ) = 1,
        \label{dGdchic}
\end{equation} 
\begin{equation}
        \left ( \frac{\partial G}{\partial \delta_w} \right ) = -\chi_c^2 \left [  1 + 2x_w c S_{cw} \chi_w \left ( \delta_w - 1\right ) \right ],
        \label{dGddelta}
\end{equation} 
\begin{equation}
        \left ( \frac{\partial H}{\partial \chi_w} \right ) = 1 + 2 x_w c \delta_w \chi_w^2,
        \label{dHdchiw}
\end{equation} 
\begin{equation}
        \left ( \frac{\partial H}{\partial \chi_c} \right ) = 2 x_c c S_{cw} \delta_w \chi_c^2,
        \label{dHdchic}
\end{equation} 
\begin{equation}
        \left ( \frac{\partial H}{\partial \delta_w} \right ) = -\chi_w^2 \left [  1 + \left ( 2x_w c \chi_w + 2x_c c S_{cw} \chi_w \right ) \left ( \delta_w - 1\right ) \right ],
        \label{dHddelta}
\end{equation} 
\begin{equation}
        V\left ( \frac{\partial G}{\partial V} \right ) =-2 x_w c S_{cw} \delta_w \chi_w^3,
        \label{VdGdV}
\end{equation} 
\begin{equation}
        V\left ( \frac{\partial H}{\partial V} \right ) =-\chi_w^2 \left ( 2 x_w c \delta_w \chi_w +  2 x_c c S_{cw} \delta_w \chi_c \right ),
        \label{VdHdV}
\end{equation} 
\begin{equation}
        V\left ( \frac{\partial \delta_w}{\partial V} \right ) =-  \frac{\delta_w \eta}{g_\eta} \frac{\mathrm{d} g_\eta} {\mathrm{d} \eta},
        \label{VddeltadV}
\end{equation} 
\begin{equation}
        N\left ( \frac{\partial G}{\partial N_w} \right ) =2  c S_{cw} \delta_w \chi_w^3,
        \label{NdGdNw}
\end{equation} 
\begin{equation}
        N\left ( \frac{\partial H}{\partial N_w} \right ) =2  c \delta_w \chi_w^3,
        \label{NdHdNw}
\end{equation} 
\begin{equation}
        N\left ( \frac{\partial H}{\partial N_c} \right ) =2  c S_{cw} \delta_w \chi_w^2 \chi_c,
        \label{NdHdNc}
\end{equation} 
\begin{equation}
        N\left ( \frac{\partial \delta_w}{\partial N_i} \right ) = \frac{\delta_w \eta b_i}{g_\eta b} \frac{\mathrm{d} g_\eta} {\mathrm{d} \eta}.
        \label{NddeltadNi}
\end{equation}


\subsection{Newton Solver for the Aqueous Inner Loop}

Computing $\ln\phi_i$ in the aqueous phase requires the simultaneous solution of
three self-consistency equations in the unknowns
$\mathbf{v} = (Z_w,\,\varepsilon_r,\,\chi_w)^\top$:
\begin{align}
  F_0(Z_w, \varepsilon_r, \chi_w) &= Z_w - Z_w^{\text{phys}} - Z_w^{\text{assoc}}
    - Z_w^{\text{DH}} - Z_w^{\text{Perm}} = 0, \label{Eq:F0} \\
  F_1(Z_w, \varepsilon_r, \chi_w) &= T_1(\varepsilon_r,\varepsilon_\infty)
    - T_2(Z_w, \chi_w) = 0, \label{Eq:F1} \\
  F_2(Z_w, \varepsilon_r, \chi_w) &= \chi_w - \chi_w^{\text{new}}(Z_w, \chi_w) = 0. \label{Eq:F2}
\end{align}
Here $Z_w^{\text{phys}}$, $Z_w^{\text{assoc}}$, $Z_w^{\text{DH}}$ and $Z_w^{\text{Perm}}$
are defined in Eqs.~\eqref{Eq:Z_phys}--\eqref{Eq:ZPerm}, $T_1$ and $T_2$ are,
respectively, the left- and right-hand sides of the relative permittivity
self-consistency equation~\eqref{Eqeps_simp},
and $\chi_w^{\text{new}}$ follows from the H-function, Eq.~\eqref{Hfunction}.

Writing the dimensionless association strength $\Delta = \delta_w P/(RT)$, the
cross-associating fraction is
\begin{equation}
  \chi_c = \frac{Z_w}{Z_w + 2x_w \chi_w S_{cw} \Delta}, \qquad
  D_c \equiv Z_w + 2x_w \chi_w S_{cw} \Delta,
  \label{Eq:chic_Zform}
\end{equation}
and the aqueous H$_2$O fraction reads
\begin{equation}
  \chi_w^{\text{new}} = \frac{Z_w}{Z_w + 2x_w \chi_w \Delta + 2x_c \chi_c S_{cw} \Delta},
  \qquad D_w \equiv Z_w + 2x_w \chi_w \Delta + 2x_c \chi_c S_{cw} \Delta.
  \label{Eq:chiw_Zform}
\end{equation}

Newton's method is applied to Eqs.~\eqref{Eq:F0}--\eqref{Eq:F2}:
\begin{equation}
  \mathbf{v}_{k+1} = \mathbf{v}_k - \mathbf{J}(\mathbf{v}_k)^{-1}\,
  \mathbf{F}(\mathbf{v}_k),
  \label{Eq:Newton}
\end{equation}
where the $3\times 3$ analytical Jacobian $J_{ij}=\partial F_i/\partial v_j$ is
derived below.

\subsubsection*{Row 2: $\partial F_2 / \partial \mathbf{v}$}

Because $F_2 = \chi_w - Z_w/D_w$, the chain of dependencies
$\chi_c(\chi_w)$ and $\chi_c(Z_w)$ must be carried through $D_w$.

The partial of $\chi_c$ with respect to $Z_w$ is
\begin{equation}
  \frac{\partial \chi_c}{\partial Z_w} = \frac{D_c - Z_w\,\partial D_c/\partial Z_w}{D_c^2},
  \qquad
  \frac{\partial D_c}{\partial Z_w} = 1 + 2x_w \chi_w S_{cw} \frac{\partial \Delta}{\partial Z_w},
  \label{Eq:dchic_dZw}
\end{equation}
with $\partial\Delta/\partial Z_w$ obtained from $\Delta = \delta_w P/(RT)$ and
$\rho \propto Z_w^{-1}$ at fixed $T$, $P$:
\begin{equation}
  Z_w \frac{\partial \Delta}{\partial Z_w}
  = Z_w \frac{\partial \delta_w}{\partial Z_w}\frac{P}{RT}
  = V\frac{\partial \delta_w}{\partial V}\frac{P}{RT}
  = -\frac{\Delta \eta}{g_\eta}\frac{\mathrm{d}g_\eta}{\mathrm{d}\eta}.
  \label{Eq:dDelta_dZw}
\end{equation}
Analogously, $\partial \chi_c/\partial\chi_w = -Z_w \cdot 2x_w S_{cw}\Delta / D_c^2$.
With $D_w$ defined in Eq.~\eqref{Eq:chiw_Zform}, the Jacobian entries are:
\begin{align}
  J_{20} &= -\frac{\partial \chi_w^{\text{new}}}{\partial Z_w}
    = -\frac{D_w - Z_w\,\partial D_w/\partial Z_w}{D_w^2}, \label{Eq:J20} \\
  J_{21} &= 0, \label{Eq:J21} \\
  J_{22} &= 1 - \frac{\partial \chi_w^{\text{new}}}{\partial \chi_w}
    = 1 + \frac{Z_w\,\partial D_w/\partial \chi_w}{D_w^2}, \label{Eq:J22}
\end{align}
where
\begin{align}
  \frac{\partial D_w}{\partial Z_w} &=
    1 + \left(2x_w\chi_w + 2x_c\chi_c S_{cw}\right)\frac{\partial \Delta}{\partial Z_w}
    + 2x_c S_{cw} \Delta\,\frac{\partial \chi_c}{\partial Z_w}, \\
  \frac{\partial D_w}{\partial \chi_w} &=
    2x_w\Delta + 2x_c S_{cw} \Delta\,\frac{\partial \chi_c}{\partial \chi_w}.
\end{align}

\subsubsection*{Row 1: $\partial F_1 / \partial \mathbf{v}$}

$T_1 = F(\varepsilon_r,\varepsilon_\infty)$ is the left-hand side of
Eq.~\eqref{Eqeps_simp}.  Its partial derivatives are those already given in
Eqs.~\eqref{dFder} and \eqref{dFdeinf}.
$T_2 = (N_A c / 9\varepsilon_0 k_B T)\,x_w g_w \mu_{0w}^2$ depends on $Z_w$ through
$\rho = P/(Z_w RT)$ and on $\chi_w$ through $g_w$.

The three entries are:
\begin{align}
  J_{10} &= \left(\frac{\partial F}{\partial\varepsilon_\infty}\right)
             \frac{V(\partial\varepsilon_\infty/\partial V)}{Z_w}
             - \frac{N_A\mu_{0w}^2}{9\varepsilon_0 k_B T}\,x_w
             \left(\frac{\partial\rho}{\partial Z_w}\,g_w
                   + \rho\,\frac{\partial g_w}{\partial Z_w}\right), \label{Eq:J10} \\
  J_{11} &= \left(\frac{\partial F}{\partial\varepsilon_r}\right)
           = \frac{2\varepsilon_r^2 + \varepsilon_\infty^2}
                  {\varepsilon_r^2(\varepsilon_\infty+2)^2}, \label{Eq:J11} \\
  J_{12} &= -\frac{N_A\rho\mu_{0w}^2}{9\varepsilon_0 k_B T}\,x_w\,
              \frac{\partial g_w}{\partial \chi_w}, \label{Eq:J12}
\end{align}
where $V(\partial\varepsilon_\infty/\partial V)$ is given by Eq.~\eqref{deinfdV},
$\partial\rho/\partial Z_w = -\rho/Z_w$, and $\partial g_w/\partial Z_w$ follows
from differentiating Eq.~\eqref{Kirkwood_simp} through $P_{ww}$ and $P_{wc}$
(Eqs.~\eqref{dPwwdV}--\eqref{dPwcdNi}) at fixed $\chi_w$, $\chi_c$, using
Eq.~\eqref{Eq:dDelta_dZw} for $\partial\delta_w/\partial Z_w$.  Similarly,
$\partial g_w/\partial\chi_w$ uses $\partial\chi_c/\partial\chi_w$ derived above.

\subsubsection*{Row 0: $\partial F_0 / \partial \mathbf{v}$}

The Debye screening length $\kappa = \left(e^2 N_A\rho\sum_j x_j z_j^2 /
k_BT\varepsilon_0\varepsilon_r\right)^{1/2}$ gives
\begin{equation}
  \frac{\partial\kappa}{\partial Z_w} = -\frac{\kappa}{2Z_w}, \qquad
  \frac{\partial\kappa}{\partial\varepsilon_r} = -\frac{\kappa}{2\varepsilon_r}.
  \label{Eq:dkappa}
\end{equation}
For the DH contribution:
\begin{equation}
  J_{01} = -\frac{\partial Z^{\text{DH}}}{\partial\varepsilon_r}
           - \frac{\partial Z^{\text{Perm}}}{\partial\varepsilon_r}.
  \label{Eq:J01}
\end{equation}
The first term follows from differentiating Eq.~\eqref{Eq:ZDH} through
$\kappa(\varepsilon_r)$ using Eq.~\eqref{Eq:dkappa}.  The second requires
differentiating $Z^{\text{Perm}} = -a_{DH+B,\varepsilon_r}\,V(\partial\varepsilon_r/\partial V)$
with respect to $\varepsilon_r$, yielding
\begin{equation}
  \frac{\partial Z^{\text{Perm}}}{\partial\varepsilon_r} =
    -\left[\frac{\partial a_{DH+B,\varepsilon_r}}{\partial\varepsilon_r}\,V\!\left(\frac{\partial\varepsilon_r}{\partial V}\right)
    + a_{DH+B,\varepsilon_r}\,\frac{\partial}{\partial\varepsilon_r}
      V\!\left(\frac{\partial\varepsilon_r}{\partial V}\right)\right],
  \label{Eq:dZPerm_der}
\end{equation}
where
\begin{equation}
  \frac{\partial a_{DH+B,\varepsilon_r}}{\partial\varepsilon_r}
  = \frac{a_{DH+B,\varepsilon_r}}{H_{DH}}\,
    \frac{\partial H_{DH}}{\partial\kappa}\,
    \frac{\partial\kappa}{\partial\varepsilon_r}
    - \frac{2\,a_{DH+B,\varepsilon_r}}{\varepsilon_r},
  \label{Eq:daDHBder}
\end{equation}
with $H_{DH} = \sum_j x_j z_j^2\!\left(\kappa/(1+\kappa\sigma_j) - 1/R_{B,j}\right)$
and $\partial H_{DH}/\partial\kappa = \sum_j x_j z_j^2/(1+\kappa\sigma_j)^2$.
The derivative of $V(\partial\varepsilon_r/\partial V)$ with respect to $\varepsilon_r$
is obtained by differentiating Eq.~\eqref{Ffunction_V} through $\partial F/\partial\varepsilon_r$
and $\partial F/\partial\varepsilon_\infty$ (which both depend on $\varepsilon_r$):
\begin{equation}
  \frac{\partial}{\partial\varepsilon_r}V\!\left(\frac{\partial\varepsilon_r}{\partial V}\right)
  = -\frac{V(\partial\varepsilon_r/\partial V)}{\partial F/\partial\varepsilon_r}\,
     \frac{\partial^2 F}{\partial\varepsilon_r^2}
  - \frac{1}{\partial F/\partial\varepsilon_r}\,
    \frac{\partial^2 F}{\partial\varepsilon_r\,\partial\varepsilon_\infty}\,
    V\!\left(\frac{\partial\varepsilon_\infty}{\partial V}\right),
  \label{Eq:d_VderdV_der}
\end{equation}
where the second-order partial derivatives of $F$ (Eq.~\eqref{Eqeps_simp}) are
\begin{equation}
  \frac{\partial^2 F}{\partial\varepsilon_r^2}
  = -\frac{2\varepsilon_\infty^2}{\varepsilon_r^3(\varepsilon_\infty+2)^2}, \qquad
  \frac{\partial^2 F}{\partial\varepsilon_r\,\partial\varepsilon_\infty}
  = \frac{4(\varepsilon_\infty - \varepsilon_r^2)}{\varepsilon_r^2(\varepsilon_\infty+2)^3}.
  \label{Eq:d2F}
\end{equation}

The entry $J_{02}$ arises from the $\chi_w$-dependence of $Z^{\text{assoc}}$ and of
$Z^{\text{Perm}}$ (through $V(\partial\varepsilon_r/\partial V)$, which depends on
$g_w$ and thus on $\chi_w$):
\begin{equation}
  J_{02} = -\frac{\partial Z^{\text{assoc}}}{\partial\chi_w}
           - \frac{\partial Z^{\text{Perm}}}{\partial\chi_w},
  \label{Eq:J02}
\end{equation}
where $\partial Z^{\text{assoc}}/\partial\chi_w$ follows directly from
Eq.~\eqref{Eq:Z_assoc} (differentiating through $\chi_c(\chi_w)$), and
$\partial Z^{\text{Perm}}/\partial\chi_w = a_{DH+B,\varepsilon_r}\,
(\partial V(\partial\varepsilon_r/\partial V)/\partial\chi_w)$ requires
differentiating the chi cross-derivative system (Eqs.~\eqref{Gfunction_Phi}
and \eqref{Hfunction_Phi}) with respect to $\chi_w$ to obtain
$\partial(V\partial\chi_w/\partial V)/\partial\chi_w$ and
$\partial(V\partial\chi_c/\partial V)/\partial\chi_w$, and then propagating
through $P_{ww}$, $P_{wc}$, $g_w$, and $V(\partial F/\partial V)$.

The entry $J_{00}$ approximates $\partial(Z^{\text{phys}}+Z^{\text{assoc}}+Z^{\text{DH}}
+Z^{\text{Perm}})/\partial Z_w$ by omitting one chain-rule contribution:
the term arising from the implicit $Z_w$-dependence of
$V(\partial\chi_w/\partial V)$ and $V(\partial\chi_c/\partial V)$ through the
chi cross-derivative system.  Validating against centered finite differences
shows this residual to be 1--2\,\% at typical conditions, which does not
impair Newton convergence from warm starts.

